\subsection{SHAP Explainability and Feature-Ranking Stability (RQ3)}
\label{subsec:shap_methodology}

To address RQ3, this study uses SHAP (SHapley Additive exPlanations) to quantify
how strongly each input feature contributes to model predictions. SHAP assigns
an additive attribution value to each feature for an individual prediction and
is grounded in Shapley-value principles from cooperative game theory
\cite{lundberg2017shap}. For tree ensembles, TreeSHAP provides an efficient
algorithm for obtaining these attributions while preserving the consistency
properties of SHAP \cite{lundberg2020treeexplainer}.

The explainability analysis follows the same stratified five-fold
cross-validation partitions used for predictive evaluation. For each dataset,
model, and fold, the model is fitted only on the training portion of that fold.
SHAP values are then computed on a deterministic sample of observations from
the corresponding validation portion. This design keeps the explained
observations separate from the rows used to fit the model and allows changes in
feature importance to be measured as the training sample changes across folds.
The random seed, maximum number of explained validation observations, and other
SHAP settings are stored in \texttt{config.yaml} so that the analysis is
reproducible.

For XGBoost and LightGBM, the implementation uses TreeSHAP through
\texttt{TreeExplainer}. The scikit-learn AdaBoost classifier used in this study
is not supported by TreeSHAP in the pinned SHAP implementation, so AdaBoost is
explained using model-agnostic permutation SHAP with a deterministic background
sample drawn from the fold-training data. This difference is treated as a
computational implementation detail: the stability analysis compares feature
rankings \emph{within the same model across folds}, rather than comparing raw
SHAP magnitudes between different explainer algorithms.

For a feature $j$ in fold $f$, global importance is defined as the mean absolute
SHAP value across the $N_f$ explained validation observations:
\begin{equation}
I_{j,f} = \frac{1}{N_f}\sum_{i=1}^{N_f} |\phi_{i,j,f}|,
\end{equation}
where $\phi_{i,j,f}$ is the SHAP value for feature $j$ and observation $i$ in
fold $f$. Features are sorted in descending order of $I_{j,f}$ to produce one
feature-importance ranking for every model and fold. Mean absolute SHAP is used
because positive and negative contributions should not cancel when estimating
global importance.

Ranking stability is evaluated in two complementary ways. First, Spearman rank
correlation is computed for every pair of cross-validation folds using the full
feature-ranking vectors. With five folds, this yields ten pairwise comparisons
for each dataset--model combination. A Spearman coefficient closer to one
indicates that the relative ordering of features changes little between folds.
Second, top-$k$ overlap is measured with the Jaccard index:
\begin{equation}
J(A,B) = \frac{|A \cap B|}{|A \cup B|},
\end{equation}
where $A$ and $B$ are the sets of the top-$k$ features from two folds. The
experiment uses $k=10$ as configured in \texttt{config.yaml}. The Spearman
measure evaluates stability over the complete ranking, whereas top-$k$ Jaccard
overlap focuses on whether the most influential features remain consistent.

The implementation stores fold-level mean absolute SHAP values and ranks,
pairwise fold-stability measurements, aggregated stability summaries, and a
consensus feature ranking. The consensus table reports each feature's average
rank, rank standard deviation, mean absolute SHAP value, and frequency of
appearing among the top-$k$ features across folds. These outputs provide the
evidence required to answer RQ3 without relying on a single fitted model or a
single data split.

The corresponding empirical stability results are presented in
Section~\ref{subsec:shap_results} and summarized in
Table~\ref{tab:shap_stability_results}. This separation keeps the present
subsection focused on how the explanations are computed, while the Results and
Discussion section evaluates how stable those explanations are in practice.

